\documentclass[12pt]{article}
\usepackage{geometry}
\geometry{margin=1in}

% ----------------------------------------------------------
%  Linux Penguin Theme — Based on Recommended Colors
% ----------------------------------------------------------

% --- COLOR SETUP ---
\usepackage{xcolor}

\definecolor{cream}{RGB}{246,242,219}
\definecolor{teamred}{RGB}{209,57,64}
\definecolor{navy}{RGB}{40,46,87}
\definecolor{softwhite}{RGB}{250,250,230}
\definecolor{accentyellow}{RGB}{240,200,60}

% Extra utility colors
\definecolor{muted}{RGB}{120,125,140}
\definecolor{darkcream}{RGB}{230,225,205}

% ----------------------------------------------------------
%  FONT SETUP (Optional but Recommended)
%  Works best with XeLaTeX or LuaLaTeX
% ----------------------------------------------------------
\usepackage{fontspec}

% Main document font (Google fonts — upload to Overleaf)
% If you don't upload fonts, comment these two lines.
\setmainfont{SourceSans3}[
  Path=./,
  ItalicFont  = *-Italic,
  BoldFont    = *-Bold
]

\newfontfamily\headingfont{ScienceGothic}[
  Path=./,
]

% ----------------------------------------------------------
%  SECTION STYLE
% ----------------------------------------------------------
\usepackage{titlesec}

% Section
\titleformat{\section}
  {\Large\headingfont\bfseries\color{navy}}
  {\thesection}
  {1em}
  {\rule{\linewidth}{1pt}\vspace{0.5em}\\}

% Subsection
\titleformat{\subsection}
  {\large\headingfont\bfseries\color{teamred}}
  {\thesubsection}
  {0.8em}
  {}

% Subsubsection
\titleformat{\subsubsection}
  {\bfseries\color{navy}}
  {\thesubsubsection}
  {0.6em}
  {}

% ----------------------------------------------------------
%  PAGE STYLE
% ----------------------------------------------------------
\usepackage{fancyhdr}
\pagestyle{fancy}

\fancyhf{} % clear all

\lhead{\textcolor{navy}{\headingfont \courseTitle}}
\rhead{\textcolor{teamred}{\thepage}}
\renewcommand{\headrule}{\color{teamred}\hrule height 1pt}

% ----------------------------------------------------------
%  HYPERLINK STYLE
% ----------------------------------------------------------
\usepackage{hyperref}

\hypersetup{
  colorlinks = true,
  linkcolor  = teamred,
  urlcolor   = navy,
  citecolor  = accentyellow
}

% ----------------------------------------------------------
%  CUSTOM BOX ENVIRONMENTS
% ----------------------------------------------------------
\usepackage{tcolorbox}
\tcbuselibrary{skins,breakable}

% Info box
\newtcolorbox{infobox}{
  enhanced,
  breakable,
  colback   = softwhite,
  colframe  = navy,
  coltext   = navy,
  boxrule   = 1pt,
  arc       = 4pt,
  left=6pt, right=6pt, top=6pt, bottom=6pt
}

% Warning box
\newtcolorbox{warningbox}{
  enhanced,
  breakable,
  colback   = accentyellow!30,
  colframe  = accentyellow!80!navy,
  coltext   = navy,
  boxrule   = 1pt,
  arc       = 4pt,
  left=6pt, right=6pt, top=6pt, bottom=6pt
}

% Highlight box
\newtcolorbox{highlightbox}{
  enhanced,
  breakable,
  colback   = teamred!10,
  colframe  = teamred,
  coltext   = navy,
  boxrule   = 1pt,
  arc       = 4pt,
  left=6pt, right=6pt, top=6pt, bottom=6pt
}

% ----------------------------------------------------------
%  ITEMIZE / ENUMERATE STYLE
% ----------------------------------------------------------
\usepackage{enumitem}

\setlist[itemize]{
  label=\textcolor{teamred}{\large\textbullet},
  leftmargin=1.2em,
  itemsep=4pt
}

\setlist[enumerate]{
  label=\textcolor{navy}{\arabic*.},
  leftmargin=1.4em,
  itemsep=4pt
}

% ----------------------------------------------------------
%  TITLE STYLE
% ----------------------------------------------------------
\usepackage{titling}

\pretitle{
  \begin{center}
  \headingfont\bfseries\LARGE\color{navy}
}
\posttitle{
  \end{center}
}

\preauthor{
  \begin{center}
  \color{teamred}
}
\postauthor{
  \end{center}
}

% ----------------------------------------------------------
%  OPTIONAL — SOFT PAGE BACKGROUND
%  Comment out if you want white pages
% ----------------------------------------------------------
\usepackage{pagecolor}
\pagecolor{cream}
\color{navy}

% ----------------------------------------------------------
% CUSTOM COMMANDS FOR TEMPLATE
% ----------------------------------------------------------
\newcommand{\courseTitle}{Linux Command Line Basics}
\newcommand{\courseDate}{30/11/2025}

% ----------------------------------------------------------
% END OF THEME
% ----------------------------------------------------------

\begin{document}

% Title with minimal space
\begin{center}
\headingfont\Huge\bfseries\color{navy}
\courseTitle\\[0.5em]
\Large\color{teamred}
\courseDate
\end{center}

\vspace{1em}
\begin{center}
\color{teamred}\rule{0.8\textwidth}{1pt}
\end{center}

\section*{1. Introduction to the Linux Command Line}

The Linux command line (often called the \textbf{terminal}, \textbf{shell}, or \textbf{CLI}) is a text-based interface that allows you to interact with your system quickly and efficiently. While graphical interfaces (GUIs) are great for daily tasks, the command line provides:

\begin{itemize}
    \item Faster automation
    \item Power-user flexibility
    \item Access to advanced tools
    \item A universal interface across all distros
\end{itemize}

Most distributions ship with \textbf{bash} as the default shell, though others (zsh, fish) also exist.

To open a terminal:

\begin{itemize}
    \item \textbf{Ubuntu / Debian}: \texttt{Ctrl + Alt + T}
    \item \textbf{Fedora}: Activities → Terminal
    \item \textbf{Any distro}: Search "Terminal"
\end{itemize}

\section*{2. System Information}

Linux provides several commands to check system details:

\subsection*{2.1 \texttt{uname}}

Shows basic system information:

\begin{highlightbox}
\texttt{\$ uname}\\
\texttt{\$ uname -a}\\
\texttt{\$ uname -r}\\
\texttt{\$ uname -m}
\end{highlightbox}

\subsection*{2.2 \texttt{hostnamectl}}

Shows OS name, kernel, hardware:

\begin{highlightbox}
\texttt{\$ hostnamectl}
\end{highlightbox}

\subsection*{2.3 \texttt{inxi}}

A more advanced tool (you may need to install it):

\begin{highlightbox}
\texttt{\$ sudo apt update \# update package index}\\
\texttt{\$ sudo apt install inxi}\\
\texttt{\$ inxi -F}
\end{highlightbox}

\begin{infobox}
\textbf{Note:} you need to update the package index before installing any package.
\end{infobox}

Displays CPU, GPU, RAM, kernel, drivers, etc.

\subsection*{2.4 \texttt{lsb\_release}}

Useful on Debian-based systems:

\begin{highlightbox}
\texttt{\$ lsb\_release -a}
\end{highlightbox}

\section*{3. Shell Variables}

The shell stores data (variables) temporarily during your session.

\subsection*{3.1 Viewing Variables}

\begin{highlightbox}
\texttt{\$ echo \$HOME}\\
\texttt{\$ echo \$USER}\\
\texttt{\$ echo \$PATH}
\end{highlightbox}

\subsection*{3.2 Creating Variables}

\begin{highlightbox}
\texttt{\$ NAME="Khaled"}\\
\texttt{\$ echo \$NAME}
\end{highlightbox}

\subsection*{3.3 Exporting Variables}

Makes them available to programs:

\begin{highlightbox}
\texttt{\$ export EDITOR=nano}
\end{highlightbox}

\subsection*{3.4 Listing All Variables}

\begin{highlightbox}
\texttt{\$ printenv}\\
\texttt{\$ set}
\end{highlightbox}

\section*{4. File System Navigation}

Every Linux filesystem starts from the \textbf{root (\texttt{/})} directory.

\subsection*{4.1 \texttt{pwd} – Print Working Directory}

Shows where you currently are:

\begin{highlightbox}
\texttt{\$ pwd}
\end{highlightbox}

\subsection*{4.2 \texttt{ls} – List Files}

\begin{highlightbox}
\texttt{\$ ls}\\
\texttt{\$ ls -l}\\
\texttt{\$ ls -la}\\
\texttt{\$ ls -h}
\end{highlightbox}

\subsection*{4.3 \texttt{cd} – Change Directory}

\begin{highlightbox}
\texttt{\$ cd /home}\\
\texttt{\$ cd Documents}\\
\texttt{\$ cd ..}\\
\texttt{\$ cd \textasciitilde}\\
\texttt{\$ cd -}
\end{highlightbox}

Shortcuts:

\begin{itemize}
    \item \texttt{.} → current directory
    \item \texttt{..} → parent directory
    \item \texttt{\textasciitilde} → home directory
    \item \texttt{-} → previous directory
\end{itemize}

\section*{5. Essential File \& Directory Commands}

\subsection*{5.1 \texttt{mkdir} – Create Directory}

\begin{highlightbox}
\texttt{\$ mkdir projects}\\
\texttt{\$ mkdir -p school/assignments}
\end{highlightbox}

\subsection*{5.2 \texttt{touch} – Create Empty Files}

\begin{highlightbox}
\texttt{\$ touch notes.txt}
\end{highlightbox}

\subsection*{5.3 \texttt{cp} – Copy}

\begin{highlightbox}
\texttt{\$ cp file.txt backup.txt}\\
\texttt{\$ cp -r folder/ folder\_backup/}
\end{highlightbox}

\subsection*{5.4 \texttt{mv} – Move/Rename}

\begin{highlightbox}
\texttt{\$ mv oldname.txt newname.txt}\\
\texttt{\$ mv file.txt /home/user/Documents/}
\end{highlightbox}

\subsection*{5.5 \texttt{rm} – Remove}

\begin{warningbox}
  Dangerous but essential
\end{warningbox}

\begin{highlightbox}
\texttt{\$ rm file.txt}\\
\texttt{\$ rm -r folder/}\\
\texttt{\$ rm -rf folder/  \# very dangerous}
\end{highlightbox}

\subsection*{5.6 \texttt{cat} – View Content}

\begin{highlightbox}
\texttt{\$ cat file.txt}
\end{highlightbox}

\subsection*{5.7 \texttt{echo} – Print Text}

\begin{highlightbox}
\texttt{\$ echo "Hello Linux!"}\\
\texttt{\$ echo \$USER}
\end{highlightbox}

\subsection*{5.8 \texttt{sleep} – Pause Execution}

\begin{highlightbox}
\texttt{\$ sleep 5}
\end{highlightbox}

\section*{6. Understanding Paths}

\subsection*{6.1 Absolute Paths}

Start from the root \texttt{/}.

Example:

\begin{highlightbox}
\texttt{/home/khaled/projects/code.py}
\end{highlightbox}

\subsection*{6.2 Relative Paths}

Start from your current location.

Example:

\begin{highlightbox}
\texttt{../images/photo.png}\\
\texttt{./run.sh}
\end{highlightbox}

\section*{7. Package Management}

Every Linux distribution uses a package manager to install software.

\subsection*{7.1 \texttt{apt} – Debian, Ubuntu, Linux Mint}

\subsubsection*{Update repositories:}

\begin{highlightbox}
\texttt{\$ sudo apt update}
\end{highlightbox}

\subsubsection*{Install:}

\begin{highlightbox}
\texttt{\$ sudo apt install gedit}
\end{highlightbox}

\subsubsection*{Remove:}

\begin{highlightbox}
\texttt{\$ sudo apt remove gedit}
\end{highlightbox}

\subsubsection*{Upgrade system:}

\begin{highlightbox}
\texttt{\$ sudo apt upgrade}
\end{highlightbox}

\subsection*{7.2 \texttt{dnf} – Fedora, RHEL 8+, CentOS Stream}

\begin{highlightbox}
\texttt{\$ sudo dnf install nano}\\
\texttt{\$ sudo dnf remove nano}\\
\texttt{\$ sudo dnf update}
\end{highlightbox}

\subsection*{7.3 \texttt{yum} – Older RHEL/CentOS}

\begin{highlightbox}
\texttt{\$ sudo yum install nano}\\
\texttt{\$ sudo yum remove nano}\\
\texttt{\$ sudo yum update}
\end{highlightbox}

\section*{8. Editors}

\subsection*{8.1 nano (Terminal Editor)}

\begin{highlightbox}
\texttt{\$ nano file.txt}
\end{highlightbox}

Controls:

\begin{itemize}
    \item Save: \textbf{Ctrl + S}
    \item Exit: \textbf{Ctrl + X}
    \item Search: \textbf{Ctrl + W}
\end{itemize}

\subsection*{8.2 gedit (GUI Editor)}

\begin{highlightbox}
\texttt{\$ gedit notes.txt \&}
\end{highlightbox}

The \texttt{\&} lets the terminal continue running.

\section*{9. Getting Help}

Linux has multiple built-in help resources.

\subsection*{9.1 \texttt{man} – Manual Pages}

\begin{highlightbox}
\texttt{\$ man ls}\\
\texttt{\$ man tar}
\end{highlightbox}

\subsection*{9.2 \texttt{info}}

\begin{highlightbox}
\texttt{\$ info coreutils}
\end{highlightbox}

\subsection*{9.3 \texttt{tldr} – Simplified Examples}

Install:

\begin{highlightbox}
\texttt{\$ sudo apt install tldr}
\end{highlightbox}

Use:

\begin{highlightbox}
\texttt{\$ tldr ls}
\end{highlightbox}

\subsection*{9.4 \texttt{--help}}

\begin{highlightbox}
\texttt{\$ ls --help}\\
\texttt{\$ tar --help}
\end{highlightbox}

\subsection*{9.5 \texttt{apropos} – Search for Commands by Description}

\begin{highlightbox}
\texttt{\$ apropos network}
\end{highlightbox}

\section*{10. Archiving \& Compression}

\subsection*{10.1 \texttt{tar} (Most Common in Linux)}

Create:

\begin{highlightbox}
\texttt{\$ tar -cvf archive.tar folder/}
\end{highlightbox}

Extract:

\begin{highlightbox}
\texttt{\$ tar -xvf archive.tar}
\end{highlightbox}

Compressed with gzip:

\begin{highlightbox}
\texttt{\$ tar -czvf archive.tar.gz folder/}\\
\texttt{\$ tar -xzvf archive.tar.gz}
\end{highlightbox}

\subsection*{10.2 \texttt{zip}}

\begin{highlightbox}
\texttt{\$ zip archive.zip file1 file2}\\
\texttt{\$ unzip archive.zip}
\end{highlightbox}

\subsection*{10.3 \texttt{7z} (7-Zip)}

Install:

\begin{highlightbox}
\texttt{\$ sudo apt install p7zip-full}
\end{highlightbox}

Compress:

\begin{highlightbox}
\texttt{\$ 7z a archive.7z folder/}
\end{highlightbox}

Extract:

\begin{highlightbox}
\texttt{\$ 7z x archive.7z}
\end{highlightbox}

\section*{11. Bonus Useful Commands}

\subsection*{11.1 \texttt{clear} – clears the terminal}

\subsection*{11.2 \texttt{history} – shows previous commands}

\subsection*{11.3 \texttt{head} / \texttt{tail} – preview files}

\begin{highlightbox}
\texttt{\$ head file.txt}\\
\texttt{\$ tail file.txt}\\
\texttt{\$ tail -f log.txt}
\end{highlightbox}

\subsection*{11.4 \texttt{less} – view files with navigation}

\begin{highlightbox}
\texttt{\$ less file.txt}
\end{highlightbox}

\subsection*{11.5 \texttt{grep} – search text}

\begin{highlightbox}
\texttt{\$ grep "error" log.txt}
\end{highlightbox}

\subsection*{11.6 \texttt{whoami} – current user}

\subsection*{11.7 \texttt{df -h} – disk usage}

\subsection*{11.8 \texttt{du -sh folder/} – folder size}

\section*{12. Final Tips for Newcomers}

\begin{itemize}
    \item Practice every day
    \item Use \texttt{tldr} for quick examples
    \item Avoid \texttt{rm -rf} unless absolutely necessary
    \item Use tab completion
    \item Use \texttt{history} to track commands
    \item Explore \texttt{/etc}, \texttt{/home}, \texttt{/usr}, \texttt{/bin} directories carefully
\end{itemize}

\section*{Practice Exercises}

\begin{enumerate}
    \item Create a directory structure:
    
    \begin{highlightbox}
    \texttt{projects/python/basics}
    \end{highlightbox}
    
    \item Create three files inside \texttt{basics}.
    
    \item Write your name into \texttt{info.txt} using \texttt{echo}.
    
    \item Compress the folder using:
    \begin{itemize}
        \item tar
        \item zip
    \end{itemize}
    
    \item Install \texttt{tldr} and test 5 commands.
    
    \item View your system kernel version and CPU info using two methods.
\end{enumerate}

% Footer
\vspace{\fill}
\begin{center}
\color{muted}
All rights reserved for \texttt{TogetherTeam} [NOT For Commercial Use] !!\\
contact the author via email \href{mailto:khaledmhfz2004@gmail.com}{\texttt{khaledmhfz2004@gmail.com}}
\end{center}
\end{document}